\documentclass[10pt,a4paper]{amsart}
\usepackage{amsmath,amssymb,amsthm,graphicx}
\usepackage{verbatim}
\newenvironment{reference}{\comment}{\endcomment}
%\newenvironment{reference}{}{}
\newenvironment{slogan}{\comment}{\endcomment}
\newenvironment{history}{\comment}{\endcomment}

% For commutative diagrams we use Xy-pic
\usepackage[all]{xy}

% We use 2cell for 2-commutative diagrams.
\xyoption{2cell}
% Theorem environments.
%
\theoremstyle{plain}
\newtheorem{theorem}[subsection]{Theorem}
\newtheorem{proposition}[subsection]{Proposition}
\newtheorem{lemma}[subsection]{Lemma}

\theoremstyle{definition}
\newtheorem{definition}[subsection]{Definition}
\newtheorem{example}[subsection]{Example}
\newtheorem{exercise}[subsection]{Exercise}
\newtheorem{situation}[subsection]{Situation}

\theoremstyle{remark}
\newtheorem{remark}[subsection]{Remark}
\newtheorem{remarks}[subsection]{Remarks}

\numberwithin{equation}{subsection}

% Macros
%
\def\lim{\mathop{\mathrm{lim}}\nolimits}
\def\colim{\mathop{\mathrm{colim}}\nolimits}
\def\Spec{\mathop{\mathrm{Spec}}}
\def\Hom{\mathop{\mathrm{Hom}}\nolimits}
\def\Ext{\mathop{\mathrm{Ext}}\nolimits}
\def\SheafHom{\mathop{\mathcal{H}\!\mathit{om}}\nolimits}
\def\SheafExt{\mathop{\mathcal{E}\!\mathit{xt}}\nolimits}
\def\Sch{\mathit{Sch}}
\def\Mor{\mathop{\mathrm{Mor}}\nolimits}
\def\Ob{\mathop{\mathrm{Ob}}\nolimits}
\def\Sh{\mathop{\mathit{Sh}}\nolimits}
\def\NL{\mathop{N\!L}\nolimits}
\def\CH{\mathop{\mathrm{CH}}\nolimits}
\def\proetale{{pro\text{-}\acute{e}tale}}
\def\etale{{\acute{e}tale}}
\def\QCoh{\mathit{QCoh}}
\def\Ker{\mathop{\mathrm{Ker}}}
\def\Im{\mathop{\mathrm{Im}}}
\def\Coker{\mathop{\mathrm{Coker}}}
\def\Coim{\mathop{\mathrm{Coim}}}

% Boxtimes
%
\DeclareMathSymbol{\boxtimes}{\mathbin}{AMSa}{"02}

%
% Macros for moduli stacks/spaces
%
\def\QCohstack{\mathcal{QC}\!\mathit{oh}}
\def\Cohstack{\mathcal{C}\!\mathit{oh}}
\def\Spacesstack{\mathcal{S}\!\mathit{paces}}
\def\Quotfunctor{\mathrm{Quot}}
\def\Hilbfunctor{\mathrm{Hilb}}
\def\Curvesstack{\mathcal{C}\!\mathit{urves}}
\def\Polarizedstack{\mathcal{P}\!\mathit{olarized}}
\def\Complexesstack{\mathcal{C}\!\mathit{omplexes}}
% \Pic is the operator that assigns to X its picard group, usage \Pic(X)
% \Picardstack_{X/B} denotes the Picard stack of X over B
% \Picardfunctor_{X/B} denotes the Picard functor of X over B
\def\Pic{\mathop{\mathrm{Pic}}\nolimits}
\def\Picardstack{\mathcal{P}\!\mathit{ic}}
\def\Picardfunctor{\mathrm{Pic}}
\def\Deformationcategory{\mathcal{D}\!\mathit{ef}}
\usepackage{cleveref}
\crefname{originallemma}{original lemma}{original lemmas}
\Crefname{originallemma}{Original lemma}{Original lemmas}
\makeatletter
\newlabel{coherent-lemma-relative-affine-cohomology}{{089W}{source}{}{sourcepointer-1}{}}
\newlabel{coherent-lemma-relative-affine-cohomology@cref}{{[originallemma][0][]089W}{[1][0][]source}{}{}{}}
\makeatother
\crefname{originalproposition}{original proposition}{original propositions}
\Crefname{originalproposition}{Original proposition}{Original propositions}
\makeatletter
\newlabel{cohomology-proposition-vanishing-Noetherian}{{02UZ}{source}{}{sourcepointer-2}{}}
\newlabel{cohomology-proposition-vanishing-Noetherian@cref}{{[originalproposition][0][]02UZ}{[1][0][]source}{}{}{}}
\makeatother
\crefname{originallemma}{original lemma}{original lemmas}
\Crefname{originallemma}{Original lemma}{Original lemmas}
\makeatletter
\newlabel{varieties-lemma-chi-tensor-finite}{{0AYT}{source}{}{sourcepointer-3}{}}
\newlabel{varieties-lemma-chi-tensor-finite@cref}{{[originallemma][0][]0AYT}{[1][0][]source}{}{}{}}
\makeatother
\title{Genus cannot increase along finite morphisms that are isomorphisms over a dense open}
\author{The Stacks Project Authors}
\date{stacks-project:curves.tex:local-snapshot-20261012; GNU FDL 1.2 or later}
\begin{document}
\maketitle
\noindent\textit{Editorial note.} All delivered mathematical text is the author's own, copied from the pinned Stacks Project chapter file sources/stacks-curves/pinned-original-source.tex; no statement or proof step was rewritten, reordered or removed.\par\medskip
\noindent\textit{Editorial note.} The counted proof is the author's complete argument: \par\medskip
\noindent\textit{Editorial note.} The label lemma-genus-goes-down names the counted statement, and the proof environment immediately following it in the source is delivered verbatim as its proof.\par\medskip
\noindent\textit{Editorial note.} The counted statement is carried by the author's own lemma environment in the chapter source.\par\medskip
\noindent\textit{Allowed standard prerequisites (editorial).}\begin{itemize}
\item Higher cohomology of zero-dimensional support sheaves vanishes. Varieties, Lemma \textbackslash{}label\{varieties-lemma-chi-tensor-finite\}. Proof lines 3855-3857: the cokernel F of O\_X -> f\_*O\_Y is supported in finitely many closed points, hence has vanishing higher cohomology.
\item Vanishing of H\textasciicircum{}2 on Noetherian schemes. Cohomology, Proposition \textbackslash{}label\{cohomology-proposition-vanishing-Noetherian\}. Proof lines 3859-3861: H\textasciicircum{}2(X, G) = 0 for the coherent sheaf G on the Noetherian scheme X.
\item Cohomology commutes with affine morphisms. Cohomology of Schemes, Lemma \textbackslash{}label\{coherent-lemma-relative-affine-cohomology\}. Proof lines 3865-3867: H\textasciicircum{}1(X, f\_*O\_Y) = H\textasciicircum{}1(Y, O\_Y) because f is finite, hence affine.
\end{itemize}\medskip
\ifdefined\hypertarget\hypertarget{sourcepointer-1}{}\fi
\noindent\textit{Source reference.} Original lemma 089W: sources/stacks-chow/tags.txt, source lines 10511--10511. Prerequisite pointer invoked at proof line(s) 3874; it names a standard fact of a sibling chapter that the counted argument uses, and every step of the counted proof built on it is retained verbatim in the delivered body.\par\medskip
\ifdefined\hypertarget\hypertarget{sourcepointer-2}{}\fi
\noindent\textit{Source reference.} Original proposition 02UZ: sources/stacks-chow/tags.txt, source lines 3750--3750. Prerequisite pointer invoked at proof line(s) 3870; it names a standard fact of a sibling chapter that the counted argument uses, and every step of the counted proof built on it is retained verbatim in the delivered body.\par\medskip
\ifdefined\hypertarget\hypertarget{sourcepointer-3}{}\fi
\noindent\textit{Source reference.} Original lemma 0AYT: sources/stacks-chow/tags.txt, source lines 13801--13801. Prerequisite pointer invoked at proof line(s) 3868; it names a standard fact of a sibling chapter that the counted argument uses, and every step of the counted proof built on it is retained verbatim in the delivered body.\par\medskip
\begin{lemma}
\label{lemma-genus-goes-down}
Let $k$ be a field. Let $X$ be a proper scheme of dimension
$\leq 1$ over $k$. Let $f : Y \to X$ be a finite morphism
such that there exists a dense open $U \subset X$ over
which $f$ is a closed immersion. Then
$$
\dim_k H^1(X, \mathcal{O}_X) \geq \dim_k H^1(Y, \mathcal{O}_Y)
$$
\end{lemma}

\begin{proof}
Consider the exact sequence
$$
0 \to \mathcal{G} \to \mathcal{O}_X \to f_*\mathcal{O}_Y \to \mathcal{F} \to 0
$$
of coherent sheaves on $X$.
By assumption $\mathcal{F}$ is supported in finitely many closed points
and hence has vanishing higher cohomology
(Varieties, Lemma \ref{varieties-lemma-chi-tensor-finite}).
On the other hand, we have $H^2(X, \mathcal{G}) = 0$ by
Cohomology, Proposition \ref{cohomology-proposition-vanishing-Noetherian}.
It follows formally that the induced map
$H^1(X, \mathcal{O}_X) \to H^1(X, f_*\mathcal{O}_Y)$
is surjective. Since $H^1(X, f_*\mathcal{O}_Y) = H^1(Y, \mathcal{O}_Y)$
(Cohomology of Schemes, Lemma \ref{coherent-lemma-relative-affine-cohomology})
we conclude the lemma holds.
\end{proof}

\end{document}
